% Suggested insertion only; the manuscript has not been edited.
\subsection{Modified Gaussian benchmark: nominal quality and local sensitivity}
We studied $\min f_1=x_1^2+x_2^2$ and $\max f_2=-[(x_1-5)^2+(x_2-5)^2]+100\exp(-[(x_1-2.5)^2+(x_2-2.5)^2]/0.01)$ on $[-5,10]^2$.
An independent reduction establishes the nominal Pareto set $(t,t)$, $0\le t\le 2.5002499781$. All indicators use $(f_1,-f_2)$ with fixed reference normalization.
Twenty paired seeds per method were run for 200 generations with population and nominal offspring sizes of 100, retaining the inherited ALSF operators, thresholds and tolerances.
Mean final hypervolumes were 0.438620, 0.435319 and 0.435314 for NSGA-II, ALSF-Resample and ALSF-Drop, respectively; corresponding mean IGD+ values were 0.002122, 0.003128 and 0.003131.
ALSF-Drop used 15243.2 evaluations on average, saving 23.78\% of the nominal 20000-evaluation schedule. There were 958 actual Drop rejections within radius 0.03 of the Gaussian center, aggregated across the 20 runs.
All final nondominated designs underwent independent uncertainty-box extrema analysis. Their mean per-run sensitive fractions were 0.8805, 0.8285 and 0.8295. These full-box assessments are distinct from archive evidence labels.
Neighborhood discovery, accurate approximation of the narrow nominal Pareto segment and demonstrated rejection must be interpreted separately. Filtering can remove nominally efficient designs; lower evaluation counts alone do not establish computational advantage. Detailed per-seed results, paired bootstrap intervals, multiplicity-adjusted comparisons, censored HV targets, timing and reference accuracy are reported in the accompanying reproducibility report. No equivalence or universal superiority is inferred.
